\documentclass[11pt]{article}
\usepackage[a4paper,margin=2cm]{geometry}
\usepackage{amsmath,amssymb}
\usepackage[european]{circuitikz}
\usepackage{tcolorbox}
\tcbuselibrary{skins,breakable}
\usepackage{enumitem}
\usepackage{xcolor}
\usepackage{parskip}
\usepackage{mdframed}
\setlist{nosep,leftmargin=1.5em,topsep=2pt}

\definecolor{cblue}{HTML}{185FA5}
\definecolor{cbluebg}{HTML}{E6F1FB}
\definecolor{cteal}{HTML}{0F6E56}
\definecolor{ctealbg}{HTML}{E1F5EE}

\newtcolorbox{concept}[1]{colback=cbluebg,colframe=cblue,fonttitle=\bfseries\color{white},
  coltitle=white,title=#1,boxrule=0.6pt,arc=2pt,left=6pt,right=6pt,top=4pt,bottom=4pt}
\newtcolorbox{answer}{colback=ctealbg,colframe=cteal,boxrule=0.6pt,arc=2pt,
  left=6pt,right=6pt,top=3pt,bottom=3pt}

\newcommand{\key}[1]{\textbf{#1}}

\begin{document}

\begin{center}
{\LARGE\bfseries Th\'evenin Equivalents \& Loading}\\[2pt]
{\large End-of-day write-up}\\[2pt]
{\small 17 June 2026 \quad$\cdot$\quad Summer self-study, Day 2 PM}
\end{center}

\vspace{4pt}

\begin{concept}{Th\'evenin in three lines}
\begin{itemize}
\item Any \key{linear two-terminal network} looks, from its terminals, like one \key{$V_{Th}$ in series with one $R_{Th}$}.
\item \key{$V_{Th}$} = the \key{open-circuit voltage} (measured with nothing connected).
\item \key{$R_{Th}$} = the resistance \key{looking back in} with sources killed: a \key{voltage source $\to$ short}, a \key{current source $\to$ open}.
\end{itemize}
\end{concept}

\section*{Worked example: a loaded voltage divider}

\textbf{Problem.} A $10\,\mathrm{k}\Omega$ / $10\,\mathrm{k}\Omega$ divider runs from $5\,$V.
\textbf{(a)} Find its Th\'evenin equivalent.
\textbf{(b)} Find the output voltage when loaded by $10\,\mathrm{k}\Omega$, $1\,\mathrm{k}\Omega$ and $100\,\Omega$.

\smallskip
\textbf{Assumptions \& laws.} Ideal $5\,$V source; resistor values exact; voltage-divider rule and Th\'evenin reduction; meter assumed not to load the circuit.

\subsection*{(a) The Th\'evenin equivalent}

\textbf{Open-circuit voltage} (no load $\Rightarrow$ no current through the divider, clean ratio):
\[
V_{Th} = 5 \times \frac{10\,\mathrm{k}}{10\,\mathrm{k}+10\,\mathrm{k}} = 2.5\,\text{V}
\]

\textbf{Look-back resistance} --- kill the source (short the $5\,$V); the two $10\,\mathrm{k}$ now hang from the output to ground in \key{parallel}:
\[
R_{Th} = 10\,\mathrm{k} \parallel 10\,\mathrm{k} = \frac{10\,\mathrm{k}\times 10\,\mathrm{k}}{10\,\mathrm{k}+10\,\mathrm{k}} = 5\,\mathrm{k}\Omega
\]

\begin{answer}
\[ \boxed{V_{Th} = 2.5\,\text{V}} \qquad \boxed{R_{Th} = 5\,\mathrm{k}\Omega} \]
\end{answer}

\begin{center}
\begin{circuitikz}[european, scale=0.92, transform shape]
\draw (0,3.4) node[above]{$5\,$V} to[R, l=$10\,\mathrm{k}\Omega$] (0,1.7)
  to[R, l_=$10\,\mathrm{k}\Omega$] (0,0) node[ground]{};
\draw (0,1.7) to[short, -o] (1.7,1.7) node[right]{$X$};
\node[circ] at (0,1.7){};
\draw[->, thick] (2.9,1.7) -- (4.2,1.7) node[midway, above]{Th\'evenin};
\draw (5.4,0) node[ground]{} to[V=$2.5\,$V] (5.4,1.7)
  to[R, l=$5\,\mathrm{k}\Omega$] (5.4,3.4) to[short, -o] (7.0,3.4) node[right]{$X$};
\end{circuitikz}
\end{center}

\subsection*{(b) Output under load}

The load $R_L$ and $R_{Th}$ form a fresh divider, so:
\[
V_{out} = V_{Th}\,\frac{R_L}{R_{Th}+R_L} = 2.5\,\frac{R_L}{5\,\mathrm{k}+R_L}
\]

\begin{itemize}
\item $R_L = 10\,\mathrm{k}\Omega:\quad 2.5\times\dfrac{10}{5+10} = \mathbf{1.67\,V}$
\item $R_L = 1\,\mathrm{k}\Omega:\quad 2.5\times\dfrac{1}{5+1} = \mathbf{0.42\,V}$
\item $R_L = 100\,\Omega:\quad 2.5\times\dfrac{0.1}{5+0.1} = \mathbf{0.049\,V}$
\end{itemize}

\begin{answer}
\textbf{Key insight.} The output is \key{not} a stiff $2.5\,$V. The load pulls current through $R_{Th}$, dropping voltage across it, so less reaches the output. The \key{smaller the load relative to $R_{Th}$, the bigger the drop} --- a $100\,\Omega$ load ($\ll 5\,\mathrm{k}$) almost collapses it to zero. Rule of thumb: keep the load $\ge 10\times R_{Th}$ to stay within $\sim$10\%.
\end{answer}

\section*{What the labs taught --- the keepers}

\textbf{Meter habits}
\begin{itemize}
\item \key{Voltage} $\to$ meter in \key{parallel} (across the part). \key{Current} $\to$ meter in \key{series} (in the loop, mA socket). Never an ammeter across the supply. \key{Resistance} $\to$ only on an \key{isolated, unpowered} part.
\item Scope ground clips are all \key{common and earthed} $\to$ use \key{one ground}. For a part not touching ground, read both ends to ground and \key{subtract} (Math A$-$B). A channel reads \key{(tip) $-$ (its own ground)}.
\item \key{Every meter loads} the circuit (its $R_{in}$ vs $R_{Th}$): high $R_{in}$ = negligible, low $R_{in}$ = big drop. The \key{$10\times$ rule}.
\end{itemize}

\textbf{Core laws}
\begin{itemize}
\item \key{KVL}: drops round a loop add up to the supply ($V_R + V_{LED} = V_s$).
\item \key{Ohm's law} $V = IR$ --- pulled from your own current sweep.
\item \key{Power} $P = VI = I^2 R = V^2/R$; every part has a \key{wattage rating} --- exceed it $\to$ heat $\to$ failure. A \key{fuse} is a sacrificial link that blows above its rated current.
\end{itemize}

\textbf{The LED}
\begin{itemize}
\item It \key{clamps its voltage at $\sim$2\,V} and does \key{not limit its own current}. What destroys it is \key{current} (heat), not voltage.
\item So it needs a \key{series resistor}: the resistor takes the leftover voltage ($V_R = V_{wiper}-2$) and \key{sets the current}, $I = (V_{wiper}-2)/470$.
\end{itemize}

\textbf{Th\'evenin \& loading}
\begin{itemize}
\item Any source/divider at two terminals $=$ \key{$V_{Th}$ in series with $R_{Th}$} (open-circuit voltage; resistance with sources killed).
\item \key{Loading drops the output}: $V = V_{Th}\,R_L/(R_{Th}+R_L)$ --- watched $2.5\,$V fall to $0.049\,$V.
\item A \key{pot is an adjustable divider}: one track split by the wiper into $R_{top}+R_{bot}$ (always $=$ full value); the wiper voltage is the output, sliding $0\to5\,$V.
\end{itemize}

\textbf{Real sources}
\begin{itemize}
\item A real battery/PSU $=$ \key{ideal source $+$ internal resistance} (a Th\'evenin source). Shorting gives a \key{finite} current ($V/R_\text{internal}$); a current-limited PSU caps it.
\end{itemize}

\end{document}
